% 3.4 The psychology of constraints
\subsection{The psychology of constraints}
\label{sec:psychology}

The cases so far are about systems: grids, banks, hospitals. But there is also
a question about people. Does having less actually change how we think about a
problem, or do constraints only change which solutions survive?

There is reasonable evidence that it changes how we think, at least within
limits. In a set of experiments, \textcite{moreau2005} asked people to design
products under different restrictions. When they were limited in which
components they could use, their designs became more creative, because the
restriction blocked the path of least resistance: retrieving a familiar example
and copying it. \textcite{mehta2016} found something similar when they made
people feel that resources were scarce: they became more inventive in finding
new uses for ordinary objects, apparently because scarcity loosened the habit
of seeing each object as fit for only one purpose.
\textcite{stokes2006} makes a historical version of the same argument: many
breakthroughs in art and science came from people who deliberately narrowed
their options and then searched hard within them.

This is the psychological core of the reframing step in
\cref{fig:mechanism}. When the familiar answer is off the table, people are
pushed to ask what the problem really is.

The same literature, though, contains its own warning. In the most thorough
review I know of, \textcite{acar2019} pulled together findings on constraints
from management, psychology, engineering and the arts, and found them pointing
in different directions: some studies show constraints helping, others hurting,
others no effect at all. Their proposed resolution is an inverted U. A little
constraint raises motivation and forces wider search; too much drains
motivation and narrows attention to survival (\cref{fig:inverted-u}). They
present this as a framework that fits the scattered evidence, not as an
estimated curve, and I draw it in the same spirit.

\begin{figure}[htbp]
\centering
% Figure: the inverted U proposed in the constraints literature. Conceptual:
% the curve is drawn, not estimated, and has no units.
\begin{tikzpicture}
\begin{axis}[
  scarcity axis,
  width=0.86\linewidth, height=0.46\linewidth,
  xmin=0, xmax=10, ymin=0, ymax=10.5,
  xtick=\empty, ytick=\empty, ymajorgrids=false,
  xlabel={Tightness of constraint \quad (slack $\rightarrow$ severe)},
  ylabel={Novel, useful output},
  axis lines=left, axis line style={draw=stone, -{Stealth[length=4pt]}}, axis on top,
]
  % Shaded zones
  \fill[indigowash] (axis cs:0,0) rectangle (axis cs:3,10.5);
  \fill[amberwash] (axis cs:3,0) rectangle (axis cs:7,10.5);
  \fill[claywash] (axis cs:7,0) rectangle (axis cs:10,10.5);
  \addplot[series primary, domain=0:10, samples=100, smooth] {1.6 + 6.2*exp(-((x-4.8)^2)/7)};

  \node[note, text=indigo, text width=26mm] at (axis cs:1.5,9.4) {\textbf{Slack}\\ the path of least resistance: copy what exists};
  \node[note, text=amber!70!ink, text width=34mm] at (axis cs:5,9.95) {\textbf{Moderate}\\ familiar paths blocked; problem reframed};
  \node[note, text=clay, text width=26mm] at (axis cs:8.5,9.4) {\textbf{Severe}\\ attention consumed by the shortfall; going without};
  \node[note, anchor=south east, font=\sffamily\scriptsize\itshape] at (axis cs:10,0.2) {Conceptual sketch, not data};
\end{axis}
\end{tikzpicture}

\caption{\takeaway{Some constraint helps; severe constraint hurts.} A sketch of
the inverted-U relationship between constraints and creativity proposed by
\textcite{acar2019}. The shape is conceptual: no study has estimated where the
peak lies for households or firms facing real deprivation.
\figsource{Drawn after \textcite{acar2019}; the severe end draws on
\textcite{mani2013} and \textcite{mullainathan2013}.}}
\label{fig:inverted-u}
\end{figure}

The severe end of that curve is not hypothetical. \textcite{mani2013} found
that simply asking people to think about a difficult financial problem lowered
the cognitive performance of poorer participants, and that the same farmers
performed better on cognitive tests after harvest, when they had money, than
before it. \textcite{mullainathan2013} describe the mechanism as tunnelling:
scarcity captures attention and leaves less ``bandwidth'' for everything else,
including planning and experimenting. Someone choosing between school fees and
medicine this week is not in a good position to redesign the roof.

\begin{counterpoint}
The laboratory studies that show constraints helping mostly use mild,
temporary and artificial constraints: a restricted set of parts, a primed
feeling of scarcity. The people in my case studies face severe, lasting and
real ones. The evidence that real deprivation helps creativity is much weaker
than the evidence that it hurts cognition. Read that way, the psychology
suggests that Moser, a mechanic with a roof, tools and some spare time, was
facing a moderate constraint rather than a severe one, and that the families
with no roof, no tools and no spare time are exactly the ones who do not show
up in stories like his.
\end{counterpoint}

I think that counterpoint is largely right, and it sharpens the thesis rather
than sinking it. The useful kind of scarcity is the absence of a \emph{system},
experienced by people who still have \emph{some} slack: enough security to
experiment, enough tools to build with. Where scarcity reaches the person
themselves, the curve turns down.
